\section{Introduction}
\label{sec:intro}

% Paragraph plan (one message each):
%   1. Problem: generators can produce unlimited labelled images; when does adding them help a
%      classifier trained on scarce real data?
%   2. What is already known (synthetic data helps few-shot; FID predicts downstream use poorly) and
%      what is open (which set helps, and whether the answer depends on initialization).
%   3. This study: a controlled benchmark, paired seeds, matched step budgets, conventional
%      augmentation as the control.
%   4. Findings, then 2-4 contribution bullets of one or two lines each.
\todo{opening paragraphs 1-3}

% Moved from the report's Introduction. Its numbers are validation scores; replace them with the
% test-set numbers once the Results section is final.
\paragraph{Contribution.}
That synthetic images help when real data is scarce is expected: with 5 images per class, Stable Diffusion
class-prompt images raise ImageNet-pretrained ResNet-50 by 11.4 and ViT-B/16 by 14.5 points, in line with
earlier work on synthetic training data~\cite{he2023synthetic}. This study asks two questions that such
gains leave open. First, \emph{which property of a generated set predicts whether it will help?}
Generator-quality measures such as FID are known to predict poorly how well a classifier trained on
generated images performs~\cite{ravuri2019classification}, also for diffusion-based
augmentation~\cite{li2026diffusion}. We test this in the scarce-data setting, with the real images always
in the training set and conventional augmentation as a control. In our benchmark the two disagree:
among our three generated sets, the ones closest to the real training images help least. LoRA-adapted pools
are closer to the real data than class-prompt pools (lower FID and higher recall), and SD img2img edits of
the real images closer still, yet with pretrained classifiers the class-prompt set gives the largest gain at
every budget, and the img2img set a smaller gain than the LoRA set in 5 of 8 cells. No set-level measure we test, FID, precision or
recall, predicts which set will help, and in a first per-class analysis no precision, recall or diversity
measure predicts which classes gain. Second, \emph{why does the effect reverse with initialization?} The same images that help pretrained
classifiers lose to RandAugment for randomly initialized ones at most budgets. A recent systematic study of
diffusion-based augmentation~\cite{li2026diffusion} starts from ImageNet weights for some datasets and trains
from scratch on others, so it cannot separate the two, and it has no RandAugment control; our design pairs
both initializations on the same data, budgets and generated images. We test candidate explanations for
this sign flip (diversity of the generated images, their distance from the real data, and how well a
classifier trained on real images recognises them) and report which quality measures, if any, predict the
downstream gain. The aim is a practical rule for when to add generated images, and a warning about the
measures commonly used to choose them.

% TODO: add the GAN set (fourth generated set) and the Stage 2 checks (ConvNeXt-T, 2:1 and 5:1,
% Mixup/CutMix, stronger augmentation policies) to the paragraph above once their results are in.

\paragraph{Contributions.}
\todo{2-4 bullets, one or two lines each, each pointing to the section with its evidence}
